% GENERATED FILE. Edit canonical lessons, then run npm run generate:books.
% canonical-source-hash: fnv1a64:7cfa5b25bab2bd36
% canonical-lessons: TA-C378-amaippu, TA-C378-cankam, TA-C378-uruppinar, TA-C378-talaivar, TA-C378-porulalar, TA-R378-review-371-378

\chapter{Organisations}
\label{ch:ta-organisations}
\hlchaptermodality{378}

\begin{chapteropening}
\textbf{By the end of this chapter:} \emph{I can talk about an organisation and the people who run it.}

Complete the last lesson of chapter 378: I can talk about an organisation and the people who run it.
\end{chapteropening}

\section[\texorpdfstring{\ta{அமைப்பு} (amaippu)}{அமைப்பு (amaippu)}]{\ta{அமைப்பு} (amaippu) — an organisation; a structure}

\label{lesson:TA-C378-amaippu}



\begin{quote}
Before the new one: say the Tamil for diversity, then the Tamil for to be divided.
\end{quote}

\subsection*{You'll want to know: \ta{அமைப்பு}}
\textbf{\ta{அமைப்பு}} — \emph{amaippu} — \textquotedblleft{}an organisation; a structure\textquotedblright{}.

\begin{quote}
\ta{இந்த} \ta{அமைப்பு} \ta{ஏழைக்} \ta{குழந்தைகளுக்கு} \ta{உதவுகிறது}.
\end{quote}

\begin{quote}
\emph{inta amaippu ēḻaik kuḻantaikaḷukku utavukiṟatu.}
\end{quote}

\begin{quote}
This organisation helps poor children.
\end{quote}

Also a structure: \ta{உடலின்} \ta{அமைப்பு}, the structure of the body.

\begin{cousinweb}
\textbf{\ta{அமைப்பு}} (\emph{amaippu}) is native Dravidian: DEDR 162 lists Tamil \emph{amai}, \textquotedblleft{}to be close, attached, joined, settled\textquotedblright{}, and \emph{amaippu}: what is put together.
\end{cousinweb}

\begin{grammarlens}[title={What you've built}]
One more word for giving reasons and opinions.
\end{grammarlens}

\subsection*{Your turn}
\begin{itemize}
\raggedright
  \item \emph{Say it:} \emph{amaippu}
  \item \emph{Say it:} \emph{amaippu}, once more
  \item \emph{Say it:} \emph{inta amaippu ēḻaik kuḻantaikaḷukku utavukiṟatu.} — This organisation helps poor children.
  \item \emph{Recall:} say the Tamil for diversity, then the Tamil for to be divided, then say \emph{amaippu} again
\end{itemize}

\begin{tcolorbox}[breakable,colback=teal!4,colframe=teal!35!black,title={Before you move on}]
What does \ta{அமைப்பு} mean? (\textquotedblleft{}An organisation; a structure\textquotedblright{}.) Say it once more.
\end{tcolorbox}
\section[\texorpdfstring{\ta{சங்கம்} (caṅkam)}{சங்கம் (caṅkam)}]{\ta{சங்கம்} (caṅkam) — an association; a union}

\label{lesson:TA-C378-cankam}



\begin{quote}
Before the new one: say the Tamil for to be divided, then the Tamil for an organisation.
\end{quote}

\subsection*{You'll want to know: \ta{சங்கம்}}
\textbf{\ta{சங்கம்}} — \emph{caṅkam} — \textquotedblleft{}an association; a union\textquotedblright{}.

\begin{quote}
\ta{என்} \ta{அப்பா} \ta{ஒரு} \ta{விவசாயிகள்} \ta{சங்கத்தில்} \ta{இருக்கிறார்}.
\end{quote}

\begin{quote}
\emph{eṉ appā oru vivacāyikaḷ caṅkattil irukkiṟār.}
\end{quote}

\begin{quote}
My father belongs to a farmers' association.
\end{quote}

\ta{சங்கத்தில்} \ta{சேர்}, to join an association.

\begin{cousinweb}
\textbf{\ta{சங்கம்}} (\emph{caṅkam}) is borrowed from Sanskrit \emph{saṃgha}, \textquotedblleft{}multitude; a society, association, community\textquotedblright{} (Monier-Williams).
\end{cousinweb}

\begin{grammarlens}[title={What you've built}]
One more word for giving reasons and opinions.
\end{grammarlens}

\subsection*{Your turn}
\begin{itemize}
\raggedright
  \item \emph{Say it:} \emph{caṅkam}
  \item \emph{Say it:} \emph{caṅkam}, once more
  \item \emph{Say it:} \emph{eṉ appā oru vivacāyikaḷ caṅkattil irukkiṟār.} — My father belongs to a farmers' association.
  \item \emph{Recall:} say the Tamil for to be divided, then the Tamil for an organisation, then say \emph{caṅkam} again
\end{itemize}

\begin{tcolorbox}[breakable,colback=teal!4,colframe=teal!35!black,title={Before you move on}]
What does \ta{சங்கம்} mean? (\textquotedblleft{}An association; a union\textquotedblright{}.) Say it once more.
\end{tcolorbox}
\section[\texorpdfstring{\ta{உறுப்பினர்} (uṟuppiṉar)}{உறுப்பினர் (uṟuppiṉar)}]{\ta{உறுப்பினர்} (uṟuppiṉar) — a member}

\label{lesson:TA-C378-uruppinar}



\begin{quote}
Before the new one: say the Tamil for an organisation, then the Tamil for an association.
\end{quote}

\subsection*{You'll want to know: \ta{உறுப்பினர்}}
\textbf{\ta{உறுப்பினர்}} — \emph{uṟuppiṉar} — \textquotedblleft{}a member\textquotedblright{}.

\begin{quote}
\ta{நான்} \ta{இந்த} \ta{நூலகத்தின்} \ta{உறுப்பினர்}.
\end{quote}

\begin{quote}
\emph{nāṉ inta nūlakattiṉ uṟuppiṉar.}
\end{quote}

\begin{quote}
I am a member of this library.
\end{quote}

A word for a limb or part + -\ta{இனர்}, a person: one part of a body of people.

\begin{cousinweb}
\textbf{\ta{உறுப்பினர்}} (\emph{uṟuppiṉar}) is native: \emph{uṟuppu}, \textquotedblleft{}body, limb\textquotedblright{} (DEDR 710), with the respectful person ending -\ta{இனர்}: one limb of a body of people, so a member.
\end{cousinweb}

\begin{grammarlens}[title={What you've built}]
One more word for giving reasons and opinions.
\end{grammarlens}

\subsection*{Your turn}
\begin{itemize}
\raggedright
  \item \emph{Say it:} \emph{uṟuppiṉar}
  \item \emph{Say it:} \emph{uṟuppiṉar}, once more
  \item \emph{Say it:} \emph{nāṉ inta nūlakattiṉ uṟuppiṉar.} — I am a member of this library.
  \item \emph{Recall:} say the Tamil for an organisation, then the Tamil for an association, then say \emph{uṟuppiṉar} again
\end{itemize}

\begin{tcolorbox}[breakable,colback=teal!4,colframe=teal!35!black,title={Before you move on}]
What does \ta{உறுப்பினர்} mean? (\textquotedblleft{}A member\textquotedblright{}.) Say it once more.
\end{tcolorbox}
\section[\texorpdfstring{\ta{தலைவர்} (talaivar)}{தலைவர் (talaivar)}]{\ta{தலைவர்} (talaivar) — a leader; a president}

\label{lesson:TA-C378-talaivar}



\begin{quote}
Before the new one: say the Tamil for an association, then the Tamil for a member.
\end{quote}

\subsection*{You'll want to know: \ta{தலைவர்}}
\textbf{\ta{தலைவர்}} — \emph{talaivar} — \textquotedblleft{}a leader; a president\textquotedblright{}.

\begin{quote}
\ta{எங்கள்} \ta{சங்கத்தின்} \ta{தலைவர்} \ta{ஒரு} \ta{பெண்}.
\end{quote}

\begin{quote}
\emph{eṅkaḷ caṅkattiṉ talaivar oru peṇ.}
\end{quote}

\begin{quote}
The president of our association is a woman.
\end{quote}

\ta{தலை}, head + -\ta{வர்}: the head of a group. \ta{தலைமை}, leadership.

\begin{cousinweb}
\textbf{\ta{தலைவர்}} (\emph{talaivar}) is native Dravidian: DEDR 3103 lists Tamil \emph{talai}, \textquotedblleft{}head\textquotedblright{}, \emph{talaimai}, \textquotedblleft{}leadership\textquotedblright{}, and \emph{talaivaṉ}, \textquotedblleft{}chief\textquotedblright{}; -\ta{வர்} is the respectful form.
\end{cousinweb}

\begin{grammarlens}[title={What you've built}]
One more word for giving reasons and opinions.
\end{grammarlens}

\subsection*{Your turn}
\begin{itemize}
\raggedright
  \item \emph{Say it:} \emph{talaivar}
  \item \emph{Say it:} \emph{talaivar}, once more
  \item \emph{Say it:} \emph{eṅkaḷ caṅkattiṉ talaivar oru peṇ.} — The president of our association is a woman.
  \item \emph{Recall:} say the Tamil for an association, then the Tamil for a member, then say \emph{talaivar} again
\end{itemize}

\begin{tcolorbox}[breakable,colback=teal!4,colframe=teal!35!black,title={Before you move on}]
What does \ta{தலைவர்} mean? (\textquotedblleft{}A leader; a president\textquotedblright{}.) Say it once more.
\end{tcolorbox}
\section[\texorpdfstring{\ta{பொருளாளர்} (poruḷāḷar)}{பொருளாளர் (poruḷāḷar)}]{\ta{பொருளாளர்} (poruḷāḷar) — a treasurer}

\label{lesson:TA-C378-porulalar}



\begin{quote}
Before the new one: say the Tamil for a member, then the Tamil for a leader.
\end{quote}

\subsection*{You'll want to know: \ta{பொருளாளர்}}
\textbf{\ta{பொருளாளர்}} — \emph{poruḷāḷar} — \textquotedblleft{}a treasurer\textquotedblright{}.

\begin{quote}
\ta{பொருளாளர்} \ta{கூட்டத்தில்} \ta{செலவுகளைப்} \ta{பற்றிப்} \ta{பேசினார்}.
\end{quote}

\begin{quote}
\emph{poruḷāḷar kūṭṭattil celavukaḷaip paṟṟip pēciṉār.}
\end{quote}

\begin{quote}
The treasurer talked about the expenses at the meeting.
\end{quote}

\ta{பொருள்}, wealth + -\ta{ஆளர்}, the one who manages.

\begin{cousinweb}
\textbf{\ta{பொருளாளர்}} (\emph{poruḷāḷar}) is native: \emph{poruḷ}, \textquotedblleft{}thing, wealth\textquotedblright{} (DEDR 4544), and \emph{āḷar}, \textquotedblleft{}those who manage\textquotedblright{}, from \emph{āḷ}, \textquotedblleft{}to rule\textquotedblright{} (DEDR 5157).
\end{cousinweb}

\begin{grammarlens}[title={What you've built}]
One more word for giving reasons and opinions.
\end{grammarlens}

\subsection*{Your turn}
\begin{itemize}
\raggedright
  \item \emph{Say it:} \emph{poruḷāḷar}
  \item \emph{Say it:} \emph{poruḷāḷar}, once more
  \item \emph{Say it:} \emph{poruḷāḷar kūṭṭattil celavukaḷaip paṟṟip pēciṉār.} — The treasurer talked about the expenses at the meeting.
  \item \emph{Recall:} say the Tamil for a member, then the Tamil for a leader, then say \emph{poruḷāḷar} again
\end{itemize}

\begin{tcolorbox}[breakable,colback=teal!4,colframe=teal!35!black,title={Before you move on}]
What does \ta{பொருளாளர்} mean? (\textquotedblleft{}A treasurer\textquotedblright{}.) Say it once more.
\end{tcolorbox}
\section[(dialogue)]{Looking back: from Equal Rights to Organisations}

\label{lesson:TA-R378-review-371-378}



\begin{quote}
Say the Tamil for equality and discrimination.
\end{quote}

\subsection*{Your turn}
Say these aloud:

\begin{itemize}
\raggedright
  \item equality, discrimination, a minority, a majority, to oppress
  \item a law, to violate, to observe, to punish, to comply
  \item politics, a minister, a prime minister, a governor, a parliament
  \item an election, a candidate, to vote, a constituency, a campaign
  \item a political party, an alliance, authority, influence, to attract
  \item to struggle, a demonstration, a rally, to demand, a petition
  \item violence, non-violence, unity, diversity, to be divided
  \item an organisation, an association, a member, a leader, a treasurer
\end{itemize}

\begin{tcolorbox}[breakable,colback=teal!4,colframe=teal!35!black,title={Before you move on}]
Equality? (\textbf{\ta{சமத்துவம்}}, \emph{camattuvam}.) Discrimination? (\textbf{\ta{பாகுபாடு}}, \emph{pākupāṭu}.) A minority? (\textbf{\ta{சிறுபான்மை}}, \emph{ciṟupāṉmai}.) A majority? (\textbf{\ta{பெரும்பான்மை}}, \emph{perumpāṉmai}.) To oppress? (\textbf{\ta{ஒடுக்கு}}, \emph{oṭukku}.) A law? (\textbf{\ta{சட்டம்}}, \emph{caṭṭam}.) To violate? (\textbf{\ta{மீறு}}, \emph{mīṟu}.) To observe? (\textbf{\ta{கடைப்பிடி}}, \emph{kaṭaippiṭi}.) To punish? (\textbf{\ta{தண்டி}}, \emph{taṇṭi}.) To comply? (\textbf{\ta{இணங்கு}}, \emph{iṇaṅku}.) Politics? (\textbf{\ta{அரசியல்}}, \emph{araciyal}.) A minister? (\textbf{\ta{அமைச்சர்}}, \emph{amaiccar}.) A prime minister? (\textbf{\ta{பிரதமர்}}, \emph{piratamar}.) A governor? (\textbf{\ta{ஆளுநர்}}, \emph{āḷunar}.) A parliament? (\textbf{\ta{நாடாளுமன்றம்}}, \emph{nāṭāḷumaṉṟam}.) An election? (\textbf{\ta{தேர்தல்}}, \emph{tērtal}.) A candidate? (\textbf{\ta{வேட்பாளர்}}, \emph{vēṭpāḷar}.) To vote? (\textbf{\ta{வாக்களி}}, \emph{vākkaḷi}.) A constituency? (\textbf{\ta{தொகுதி}}, \emph{tokuti}.) A campaign? (\textbf{\ta{பிரச்சாரம்}}, \emph{piraccāram}.) A political party? (\textbf{\ta{கட்சி}}, \emph{kaṭci}.) An alliance? (\textbf{\ta{கூட்டணி}}, \emph{kūṭṭaṇi}.) Authority? (\textbf{\ta{அதிகாரம்}}, \emph{atikāram}.) Influence? (\textbf{\ta{செல்வாக்கு}}, \emph{celvākku}.) To attract? (\textbf{\ta{கவர்}}, \emph{kavar}.) To struggle? (\textbf{\ta{போராடு}}, \emph{pōrāṭu}.) A demonstration? (\textbf{\ta{ஆர்ப்பாட்டம்}}, \emph{ārppāṭṭam}.) A rally? (\textbf{\ta{பேரணி}}, \emph{pēraṇi}.) To demand? (\textbf{\ta{கோரு}}, \emph{kōru}.) A petition? (\textbf{\ta{மனு}}, \emph{maṉu}.) Violence? (\textbf{\ta{வன்முறை}}, \emph{vaṉmuṟai}.) Non-violence? (\textbf{\ta{அகிம்சை}}, \emph{akimcai}.) Unity? (\textbf{\ta{ஒற்றுமை}}, \emph{oṟṟumai}.) Diversity? (\textbf{\ta{வேற்றுமை}}, \emph{vēṟṟumai}.) To be divided? (\textbf{\ta{பிளவுபடு}}, \emph{piḷavupaṭu}.) An organisation? (\textbf{\ta{அமைப்பு}}, \emph{amaippu}.) An association? (\textbf{\ta{சங்கம்}}, \emph{caṅkam}.) A member? (\textbf{\ta{உறுப்பினர்}}, \emph{uṟuppiṉar}.) A leader? (\textbf{\ta{தலைவர்}}, \emph{talaivar}.) A treasurer? (\textbf{\ta{பொருளாளர்}}, \emph{poruḷāḷar}.)
\end{tcolorbox}
